%----------------------------------------------------------------------------------------
% ZO WERK JE MET DIT BOEK: komt vooraan in elk handboek. De uitleg over hoe de onderdelen
% werken staat alleen hier, niet meer in de hoofdstukken.
%----------------------------------------------------------------------------------------
\chapter*{Zo werk je met dit boek}

\ifhandboek
Dit boek gebruik je in de les. Je vindt er de belangrijkste leerstof, uitgewerkte voorbeelden en vooral veel opgaven. De oplossing van een opgave noteer je zelf in je schrift.
\fi
\ifzelfstudie
Deze versie is voor thuis: om een les in te halen of iets rustig te herhalen. Ze bevat hetzelfde als het handboek, met extra uitleg, extra opgaven en de antwoorden. De oefeningen hebben dezelfde nummers als in het handboek.
\fi
\ifleerkracht
Dit is de leerkrachtenversie: het handboek met didactische tips en de antwoorden. De oefeningen hebben dezelfde nummers als in het handboek.
\fi

\section*{Kleuren}

Elk hoofdstuk heeft een eigen kleur en een nummer aan de rand van de pagina. De \textbf{donkere} tint gebruiken we voor uitleg, de \textbf{felle} tint voor opdrachten.

\section*{De onderdelen}

{\renewcommand{\arraystretch}{1.75}
\begin{tabular}{@{}>{\RaggedRight\arraybackslash}p{3.3cm}>{\RaggedRight\arraybackslash}p{12.3cm}@{}}
\tikz[baseline=(k.base)]\node[fill=accent!9!white,draw=accent,line width=1pt,inner xsep=2.2mm,inner ysep=1mm,font=\sffamily\bfseries\small,text=accent](k){Theorie}; & Het omkaderde blok: definities, eigenschappen en formules, met hun bewijs. Dit moet je kennen, en hier zoek je het terug. \\
\tikz[baseline=(k.base)]{\node[inner xsep=2.2mm,inner ysep=1mm,font=\sffamily\bfseries\small,text=accent](k){Werkwijze};\draw[accent,line width=2.5pt] (k.north west) -- (k.south west);} & Het blok met een balk links: de stappen om een bepaald soort opgave aan te pakken. \\
\voorbeeldpil{Voorbeeld 1} & Een uitgewerkt voorbeeld: zo pak je dit soort opgave aan. \\
\oefstrookje{Start} & \textbf{Opwarmer:} leerstof die je nodig hebt, uit het hoofd. \textbf{Verken:} probeer dit zelf vóór de uitleg; fouten horen erbij. \\
\oefnummer{8} & Een oefening. Na elk stuk leerstof staat een korte oefening om het meteen te proberen. \\
\oefstrookje{Opdrachten} & Het grootste deel van de opgaven, na alle leerstof van het hoofdstuk. Ze staan bewust door elkaar: welke methode je nodig hebt, moet je zelf herkennen. Dit deel start op een nieuwe pagina, en vanaf daar is de tab aan de rand fel. \\
\niveau{1}\ \niveau{2}\ \niveau{3} & De moeilijkheid van een opdracht: basis, gevorderd of uitdagend. \\
\oefstrookje{Koppel} & Een kleine reeks die dit hoofdstuk verbindt met andere leerstof. De kleur van het label toont de rubriek: \koppel{analyse}{Analyse}\koppel{algebra}{Algebra}\koppel{meetkunde}{Meetkunde}\koppel{discreet}{Discrete wiskunde}\koppel{data}{Data en onzekerheid} \\
\oefstrookje{Test jezelf} & Zonder boek of notities. Verbeter daarna met de oplossingen. \\
\blokkop{Fout? Herhaal} & Dit schema staat bij de test en bij de antwoorden. Het toont wat je moet herhalen bij een fout: \oefnummer{45} $\rightarrow$ \voorbeeldpil{Voorbeeld 1} \oefnummer{9}
\end{tabular}}

\vfill
{\footnotesize\color{grijs} De kleine prenten bij de opdrachten komen uit het publiek domein, via Wikimedia Commons: tekeningen van Pearson Scott Foresman, een gravure uit \textit{Popular Science Monthly} en een houtsnede uit de 16de eeuw.\par}
